\section{Discussion}\label{sec:discussion}

\subsection{A common functional response need not imply a common mechanism}\label{sec:discussion-mechanisms}

The central empirical result is a failed mechanism-transfer hypothesis. In GPT-2 Small on IOI, L10H2 becomes more causally useful after the canonical primary Name Movers are removed, and its FALI score decomposition is target-positive and rival-negative. In Pythia-160M induction, L8H2 also becomes more causally useful after an upstream induction head is removed, but the preregistered fine-score directions reverse: the target contribution is negative and the rival contribution positive. Conditional recruitment therefore occurs in both studied settings, while the specific GPT-2 fine-score organization does not transfer.

This distinction matters because the high-level label \emph{recruitment} specifies a change in causal usefulness, not the internal route by which that change arises. Conditional-ablation methods can identify components whose importance increases after another component is removed \citep{gong2026coax}; interaction analyses show that measured component effects can depend strongly on the state of the rest of the model \citep{vaidyanathan2026curse}. Our results show why an additional mechanism-resolution step can be informative: two cases that satisfy the same functional criterion can receive opposite fine-grained score attributions under the same FALI intervention semantics.

The GPT-2 results also caution against treating a visually plausible attention adaptation as evidence of functional compensation. L10H10 increases target attention and has the same target-positive/rival-negative score pattern as L10H2, yet its total conditional gain is negative on the broad canonical distribution. Conversely, a positive or negative FALI channel should not be read in isolation as a statement that the head was recruited. Recruitment is defined by the output-grounded candidate effect; the channel values explain the endpoint change under the specified game. This separation between observed attention, component usefulness, and mechanism attribution is central to the interpretation of FALI.

The cross-setting reversal is consistent with a broader picture in which transformer mechanisms can depend on prompts, components, and model state rather than forming a single invariant circuit description \citep{franco2026promptspecific,nam2025causalhead}. Our evidence is much narrower: it concerns one GPT-2 IOI setting and one Pythia induction setting. It does not establish that fine mechanisms generally reverse across models, nor does it rule out shared mechanisms in other cases. It establishes that conditional recruitment itself is insufficient evidence for one shared target/rival mechanism.

\subsection{Relationship to conditional attribution and patching}\label{sec:discussion-positioning}

FALI is best viewed as downstream of localization and conditional-importance analysis. Activation and path patching test whether particular activations or paths causally matter for behavior \citep{goldowskydill2023pathpatching,zhang2024activationpatching,heimersheim2024activationpatching}. Conditional co-ablation asks which components become more important after a primary set is perturbed \citep{gong2026coax}. Interaction-aware mediation analysis shows that a component's causal effect can depend on the surrounding network \citep{vaidyanathan2026curse}. FALI begins after such a candidate and task-relevant target have been specified and asks a different question: which changes in the selected attention head's score computation, value content, and surrounding context account for the cross-state change in its output-grounded usefulness under an explicit counterfactual intervention model?

This role also distinguishes FALI from generic QK/OV or routing/value decompositions. The separation between attention routing and transported content is established transformer-circuit structure \citep{elhage2021mathematical}, and recent causal analyses resolve query--key routing and value-mediated computation in other settings \citep{franco2025attentioncausal,wang2026fewshot}. FALI uses these architectural objects to decompose a different estimand: the \emph{change} in the selected head's downstream causal contribution across two circuit states. The additional context factor makes explicit that an otherwise similar candidate computation can acquire different usefulness when the rest of the model changes.

The use of standard Owen attribution should likewise be understood as allocation machinery rather than a claim about unique causal pathways. Interactions are distributed according to the prospectively fixed coalition structure, and the resulting $\phi_i$ values reconcile exactly with the coarse SCORE/VALUE/CONTEXT quotient game. This guarantees accounting consistency under the defined characteristic function; it does not make the decomposition invariant to every alternative intervention or coalition design. As with other intervention-based interpretability methods, conclusions can depend on the intervention semantics \citep{miller2024faithfulness}.

These distinctions suggest a practical workflow rather than a replacement for existing tools: first localize or identify a conditionally important component, then use FALI when the scientific question is why that selected attention head's causal usefulness changed. The method is therefore complementary to circuit discovery, patching, and conditional attribution rather than a substitute for them.

\section{Limitations}\label{sec:limitations}

\paragraph{Hypothesis-driven target selection.}
FALI requires a selected attention head and a task-relevant target edge or edge set. In IOI and repeated-token induction, the target coordinates have clear task-specific interpretations. In less structured tasks, selecting a meaningful target may itself require substantial prior circuit analysis. FALI is consequently a mechanism-resolution framework, not an automatic circuit-discovery procedure.

\paragraph{Counterfactual and potentially off-manifold hybrids.}
Mixed $T/R/O/V/C$ cells combine internal states that the unmodified model need not naturally realize. Their interpretation is therefore explicitly counterfactual: they answer what the specified model would do under the intervention construction, not which latent state the model would have visited naturally. Representation-level interventions can produce internal states that diverge from natural model trajectories \citep{grant2025divergent}. We therefore do not interpret FALI attributions as proof of five independent natural causal pathways.

\paragraph{Dependence on intervention and coalition semantics.}
The decomposition depends on the intact and perturbed endpoints, the candidate-zero baseline, the structural score partition, and the fixed SCORE/VALUE/CONTEXT coalition structure. Standard Owen attribution allocates interactions consistently within that game, but another scientifically defensible intervention family or grouping could allocate interaction credit differently. Our robustness evidence is correspondingly bounded: the preregistered target-positive/rival-negative GPT-2 channels persisted under the matched-mean and matched-resample alternatives tested in EF-5, but the overall EF-5 experiment gate failed because its zero-damage reference criterion was not satisfied. We therefore do not claim robustness to arbitrary ablation choices.

\paragraph{The context factor is coarse.}
$C$ denotes the surrounding computation outside the explicitly injected candidate score/value state. It can aggregate changes in other attention heads, MLPs, residual representations, normalization, and upstream or downstream computations. A large $\phi_C$ therefore localizes explanatory credit only to the surrounding-context block of the specified game; it does not identify which external component or pathway generated that contribution. Further decomposition of $C$ would require additional hypotheses and interventions.

\paragraph{Empirical scope.}
The substantive evidence comprises two model/task settings: GPT-2 Small on IOI and Pythia-160M on repeated-token induction. The GPT-2 analysis spans broad canonical templates, multiple prompt-generation seeds, and tested ablation variants, while the Pythia test uses a separately frozen discovery/confirmation procedure. These experiments are sufficient to falsify the preregistered claim that the GPT-2 target/rival mechanism transfers unchanged to the tested Pythia induction setting, but they are not sufficient to characterize the distribution of recruitment mechanisms across transformer architectures, scales, tasks, or training stages.

\paragraph{Attention-head-specific scope.}
The current intervention algebra is designed around attention scores, values, and head outputs. Extending the same explanatory objective to MLPs, residual-stream features, or distributed multi-component mechanisms would require different architecture-aligned factors rather than mechanically reusing $T/R/O/V/C$. Thus the present results support FALI for selected attention-head mechanisms, not a universal component-attribution framework.

\section{Conclusion}\label{sec:conclusion}

We introduced FALI, an output-grounded counterfactual framework for explaining how a selected attention head's causal usefulness changes after a circuit perturbation. FALI separates target, rival, and other attention-score changes from value content and surrounding-network context, and uses standard Owen attribution to reconcile the fine score decomposition with a coarse score/value/context explanation.

The experiments establish two complementary points. First, attention adaptation and functional recruitment are distinct: a head can increase task-target attention without becoming more causally useful. Second, the same high-level functional response can conceal different mechanisms. GPT-2 IOI and Pythia induction both contain positive conditional recruitment in the studied cases, but their preregistered target/rival score contributions have opposite signs. FALI therefore provides a mechanism-discrimination layer between identifying conditional importance and asserting an internal explanation for it.

The appropriate conclusion is not that transformers share one repair mechanism, nor that FALI recovers a unique natural causal decomposition. Rather, under explicit intervention semantics, FALI makes competing mechanistic explanations of conditional attention-head importance measurable and testable. In the settings studied here, that resolution is necessary to see that a common functional label does not imply a common fine-grained mechanism.
